\subsection{Evolución del almacenamiento de credenciales}

En el corte anterior, la API de FitPulse almacenaba las contraseñas en texto plano y las comparaba directamente al iniciar sesión. Para corregir esa limitación se incorporó Argon2id en el servidor, con 19\,456 KiB de memoria, dos iteraciones y paralelismo de uno. Esta elección sigue una configuración mínima recomendada por OWASP para encarecer los intentos de adivinar contraseñas después de una filtración de la base de datos \cite{owaspPasswordStorage}. La biblioteca genera un salt aleatorio de 16 bytes en cada operación y conserva el algoritmo, los parámetros, el salt y el resultado en una cadena codificada \cite{nodeArgon2}. Por ello, la solución reutiliza la columna de contraseña existente sin agregar tablas.

El registro transforma la contraseña antes de persistirla, mientras que el inicio de sesión verifica el valor recibido contra el hash almacenado. Se mantienen los roles, la emisión de tokens y el bloqueo temporal, pero no se conserva una comparación alternativa con texto plano. Esta diferencia es una evolución del diseño del corte 2 y no una propiedad que hubiera estado implementada desde el inicio.

Para las cuentas existentes se añadió un comando de migración que actualiza únicamente la contraseña dentro de una transacción. Omite los valores que ya tienen formato Argon2id, conserva los demás datos y revierte los cambios si ocurre un fallo. Las pruebas con MySQL comprueban el acceso con la contraseña original después de migrar, la conservación de los metadatos y la repetición del procedimiento sin volver a transformar los hashes. Esta validación usa datos ficticios en una base de pruebas; cada instalación existente debe ejecutar también su propia migración antes de utilizar el login actualizado.

\subsection{Limitación de seguridad y mejora propuesta}

El hashing reduce la exposición de las contraseñas almacenadas, pero no evita el acceso mediante una contraseña correcta obtenida por phishing. En FitPulse, el impacto es mayor en las cuentas administrativas por sus permisos de modificación y consulta de ventas. Se propone exigirles un segundo factor con WebAuthn, que vincule la autenticación con el origen legítimo y reduzca la dependencia de una contraseña reutilizable \cite{owaspMultifactor}. Esta mejora no está implementada: requiere enrolamiento, desafíos, recuperación de acceso y pruebas con dispositivos compatibles, un flujo adicional al alcance de migración, registro y login abordado en este aporte.
